\section{Diverses}

\subsection{'Control'-Fragen}

\subsubsection{Wieso Feedback?}

Closed-Loop Control (Regelung) statt Open-Loop Control (Steuerung)


\subsubsection{Vorteile / Nachteile einer Steuerung}

\begin{itemize}
	\item[-] Störung $Z$ muss bekannt sein
	\item[-] Anfangswert $x_0$ muss bekannt sein
	\item[-] Struktur und Parameter des Modells müssen bekannt sein
	\item[-] (Steuer-) Strecke muss stabil sein
	\item[+]    Braucht keinen Sensor 
\end{itemize}


\subsubsection{Wozu ein Modell erstellen?}

 \begin{itemize}
	\item Steuerung: Braucht sowieso ein Modell
	\item Regelung: Modellierung von Sensor \textbf{muss} gut sein
	\item Falsches Modell kann Regler destabilisieren (Schwingung)
 \end{itemize}


\subsubsection{Vorteile linearer Modelle?}

\begin{itemize}
	\item Für lineare Modelle (lineare DGL) gibt es Standardverfahren zur Lösung 
	\item Die Systeme sind (grundsätzlich immer) lösbar
\end{itemize}



\section{Physik Formeln}

Um die DGL eines Systems aufzustellen könnten die folgenden Formeln aus der Physik hilfreich sein.


\subsection{Statik}

\subsubsection{Schiefe Ebene}

\input{tikz/schiefe_ebene.tex}

\subsubsection*{Formeln und Zusammenhänge zur schiefen Ebene}

\begin{tabular}{llll}
	$F = m \cdot a$	& $F_G = m \cdot g$ &	$F_N = m \cdot g \cdot \cos(\alpha)$ &	$F_H = m \cdot g \cdot \sin(\alpha)$
\end{tabular}
	
	
\subsubsection{Kräftegleichgewicht -- Beispiel}

\begin{minipage}[c]{0.35\columnwidth}
	\begin{align*}
		F_{\rm res} &= F_G + F_R \\
		m \cdot a	&= m \cdot g + F_N \cdot \mu
	\end{align*}
\end{minipage}
\hfill
\begin{minipage}[c]{0.6\columnwidth}
	\textbf{ \textrightarrow\ Reibung ist proportional zur Geschwindigkeit!}
\end{minipage}


\subsection{Kinematik}

\renewcommand{\arraystretch}{1.1}
\begin{tabular}{ll}
	Ort 			& $x$ 						\\
	Geschwindigkeit & $v = \dot{x}$ 			\\
	Beschleunigung 	& $a = \dot{v} = \ddot{x}$
\end{tabular}
\renewcommand{\arraystretch}{1}


\subsubsection{Gleichförmige Bewegung $\bm{a(t) = 0}$}

\renewcommand{\arraystretch}{1.3}
\begin{tabular}{l}
	$a(t) = \ddot{x} = 0$ 			\\
	$v(t) = v_0 = \dot{x} = \const$ \\
	$x(t) = v_0 \cdot t + x_0$
\end{tabular}
\renewcommand{\arraystretch}{1}
		
		
\subsubsection{Gleichmässig beschleunigte Bewegung $\bm{a(t) = \const}$}

\renewcommand{\arraystretch}{1.3}
\begin{tabular}{l}
	$a(t) = a_0 = \ddot{x} = \const$ 							\\
	$v(t) =  \dot{x} = a_0 \cdot t + v_0$ 						\\
	$x(t) = \frac{1}{2} \, a_0 \cdot t^2 + v_0 \cdot t + x_0$
\end{tabular}
\renewcommand{\arraystretch}{1}
	

\subsection{Dynamik}

\subsubsection{Translation vs. Rotation}
\renewcommand{\arraystretch}{1.2}
\begin{tabular}{| c | c |}
	\hline
	\textbf{Translation} 	& \textbf{Rotation}	\\
	\hline
	$a$ 					& $\alpha$ 			\\
	\hline
	$v$ 					& $\omega$ 			\\
	\hline
	$x$ 					& $\theta$ 			\\
	\hline
	$F$ 					& $M$ 				\\
	\hline
	$m$ 					& $J$ 				\\
	\hline
\end{tabular}
\renewcommand{\arraystretch}{1}


\subsubsection{Leistung}

\vspace{-0.2cm}
$$ \boxed{ P = \dot{W} = F \cdot v } $$


\subsubsection{Arbeit / Energie}

\vspace{-0.2cm}
$$ \boxed{ W = F \cdot s } $$

\subsubsection{Energieerhaltung}

\vspace{-0.2cm}
$$ \boxed{ W = \underbrace{m \cdot g \cdot h}_{\substack{\text{pot. Energie}}} =  m \cdot g \cdot \underbrace{ \frac{1}{2} \, g \cdot t^2}_{\substack{\text{h(t)}}}  = \underbrace{ \frac{1}{2} m \cdot v^2 }_{\substack{\text{kin. Energie}}} } $$ 
\vspace{0.2cm}


\myul{\textbf{Energiesatz der Mechanik}}
$$ \boxed{ E_{\rm pot} + E_{\rm kin} = E_{\rm tot} = \const} \qquad \text{(gilt zu jedem Zeitpunkt)} $$


\subsection{Hydrodynamik}

\subsubsection{Induzierter Widerstand $F_W$}

Induzierter Widerstand = Luftwiderstand 

$$ \boxed{ F_W = c^*_W \cdot \frac{1}{2} \cdot \rho \cdot v^2 \cdot A_{\|} } $$

\renewcommand{\arraystretch}{1.3}
\begin{tabular}{c l c}
		$F_W$ 		& Induzierter Widerstand 								& $[F_W] = \newton$ 						\\
		$c^*_W$ 	& Widerstands-Koeffizient 								& $[c*_W] = 1$ 								\\
		$\rho$ 		& Luft-Dichte		 									& $[\rho] = \frac{\kilo \gram}{\meter^3}$ 	\\
		$v$ 		& Strömungsgeschwindigkeit 								& $[v] = \frac{\meter}{\second}$ 			\\
		$A_{\|}$ 	& Projizierte Fläche \textbf{parallel} zur Strömung 	& $[A_{\|}] = \meter^2$
\end{tabular}
\renewcommand{\arraystretch}{1}


\subsection{Thermodynamik}

\subsubsection{Wärme}

\vspace{-0.2cm}
$$ \boxed{ \Delta P = \Delta \dot{Q}} $$


\subsubsection{Wärmekapazität}

\vspace{-0.2cm}
$$ \boxed{ \Delta Q = c \cdot m \cdot \Delta T = n \cdot C_m \cdot \Delta T = C \cdot \Delta T }$$


\subsubsection{Absolute Wärmekapazität $C$}

Energiespeicher-Vermögen eines \textbf{Gegenstands}
$$ \boxed{ \Delta Q = C \cdot \Delta T } $$


\subsubsection{Spezifische Wärmekapazität $c$}

Energiespeicher-Vermögen einer \textbf{Substanz}
$$ \boxed{ \Delta Q = c \cdot m \cdot \Delta T  } \qquad \qquad c_{\rm Wasser} = 4187 \, \frac{\joule}{\kilo \gram \cdot \kelvin} $$


\subsubsection{Molare Wärmekapazität $C_m$}

Energiespeicher-Vermögen einer \textbf{Anzahl Moleküle}
$$ \boxed{ C_m = \frac{c}{n} = M \cdot c } $$

\renewcommand{\arraystretch}{1.3}
\begin{tabular}{c l c}
	$\Delta Q$ 	& Zu-/Abgeführte Wärme 			& $[\Delta Q] = \joule$ 							\\
	$c$ 		& Spezifische Wärmekapazität 	& $[c] = \frac{\joule}{\kilo \gram \cdot \kelvin}$ 	\\
	$C$ 		& Absolute Wärmekapazität 		& $[C] = \frac{\joule}{\kelvin}$ 					\\
	$C_m$ 		& Molare Wärmekapazität 		& $[C_m] = \frac{\joule}{\mole \cdot \kelvin}$ 		\\
	$m$ 		& Masse 						& $[m] = \kilo \gram$ 								\\
	$\Delta T$ 	& Temperaturänderung 			& $[\Delta T] = \kelvin$ 							\\
	$n$ 		& Mol-Zahl 						& $[n] = \mole$ 									\\
	$M$ 		& Mol-Masse 					& $[M] = \frac{\kilo \gram}{\mole}$
\end{tabular}
\renewcommand{\arraystretch}{1}

